\documentclass[11pt,a4paper]{article}

\usepackage{fontspec}
\setmainfont{DejaVu Serif}[
  BoldFont={DejaVu Serif Bold},
  ItalicFont={DejaVu Serif},
  ItalicFeatures={FakeSlant=0.18},
  BoldItalicFont={DejaVu Serif Bold},
  BoldItalicFeatures={FakeSlant=0.18}
]
\setsansfont{DejaVu Sans}
\usepackage{polyglossia}
\setmainlanguage{english}
\usepackage[a4paper,margin=20.5mm]{geometry}
\usepackage{microtype}
\usepackage{setspace}
\setstretch{1.035}
\usepackage{amsmath,amssymb,amsthm,mathtools,bm,mathrsfs}
\usepackage{booktabs,longtable,array,tabularx}
\usepackage{enumitem}
\usepackage{xcolor}
\usepackage{tikz}
\usetikzlibrary{arrows.meta,positioning,calc}
\usepackage[unicode,hidelinks]{hyperref}
\usepackage[nameinlink,noabbrev]{cleveref}
\emergencystretch=3em
\allowdisplaybreaks

\hypersetup{
  pdftitle={Planck-to-Cosmos Observation Rulers},
  pdfauthor={Stassis Research Program},
  pdfsubject={Typed logarithmic scale coordinates, protocol separation, Planck normalizations, and cosmological landmarks for Geometry of Observation},
  pdfkeywords={Planck units, logarithmic coordinates, dimensional analysis, observation protocol, cosmological distance, scale geometry}
}

\definecolor{obsblue}{RGB}{37,83,138}
\definecolor{obsred}{RGB}{150,45,45}
\definecolor{obsgreen}{RGB}{30,115,78}
\definecolor{lightblue}{RGB}{239,246,253}
\definecolor{lightred}{RGB}{253,242,242}
\definecolor{lightgreen}{RGB}{239,249,244}

\theoremstyle{definition}
\newtheorem{definition}{Definition}[section]
\newtheorem{model}[definition]{Model}
\theoremstyle{plain}
\newtheorem{theorem}[definition]{Theorem}
\newtheorem{proposition}[definition]{Proposition}
\newtheorem{lemma}[definition]{Lemma}
\newtheorem{corollary}[definition]{Corollary}
\theoremstyle{remark}
\newtheorem{remark}[definition]{Remark}
\newtheorem{warning}[definition]{Boundary}

\newcommand{\R}{\mathbb{R}}
\newcommand{\Rpos}{\mathbb{R}_{>0}}
\newcommand{\X}{\mathcal{X}}
\newcommand{\Y}{\mathcal{Y}}
\newcommand{\Ocal}{\mathcal{O}}
\newcommand{\Lcal}{\mathcal{L}}
\newcommand{\Scal}{\mathcal{S}}
\newcommand{\Tcal}{\mathcal{T}}
\newcommand{\dd}{\mathop{}\!d}
\newcommand{\diag}{\operatorname{diag}}
\newcommand{\Cov}{\operatorname{Cov}}
\newcommand{\Law}{\operatorname{Law}}
\newcommand{\abs}[1]{\left\lvert #1\right\rvert}
\newcommand{\norm}[1]{\left\lVert #1\right\rVert}
\newcommand{\set}[1]{\left\{#1\right\}}
\newcommand{\status}[1]{\textsf{\footnotesize #1}}
\newcommand{\chiL}{\chi_L}
\newcommand{\chiM}{\chi_M}
\newcommand{\chiT}{\chi_T}

\newcolumntype{L}[1]{>{\raggedright\arraybackslash}p{#1}}
\newcolumntype{Y}{>{\raggedright\arraybackslash}X}

\title{\bfseries Planck-to-Cosmos Observation Rulers\\
\large Typed logarithmic charts, scale passports, and cosmological landmarks}
\author{Stassis Research Program\\
\small Independent research program; ORCID: 0009-0000-2294-705X}
\date{Strict GO Core reference revision v1.1 --- 2026}

\begin{document}
\maketitle

\begin{center}
\fcolorbox{obsblue}{lightblue}{%
\parbox{0.92\linewidth}{\small
\textbf{Status.} This document supersedes \emph{Planck-to-Cosmos
Observation Rulers v1.0}. It is a reference module for GO Core v0.2 and
the Distance-Scale Interface v0.2. The revision separates intrinsic
system scales from instrument resolutions, treats Planck units as
normalizations rather than cutoffs, and treats cosmic entries as
model-, epoch-, and convention-dependent landmarks. The logarithmic
chart is standard mathematics; no new fundamental law is claimed.}}
\end{center}

\begin{abstract}
For a positive length, invariant mass, and duration, a declared reference
triple $\bm a=(a_L,a_M,a_T)$ and base $b>1$ define the logarithmic chart
\[
 \bm\chi_{b,\bm a}(L,M,T)
 =
 \left(
 \log_b\frac{L}{a_L},
 \log_b\frac{M}{a_M},
 \log_b\frac{T}{a_T}
 \right).
\]
This map is a global smooth chart from the positive scale cone
$(\Rpos)^3$ to $\R^3$. Coherent unit changes leave it invariant, anchor
changes translate it, and base changes act by a scalar linear map.
Coordinate differences support a family of weighted log-scale metrics;
there is no canonical physical reason to assign equal weight to one
decade of length, mass, and time. Intrinsic scale catalogues belong to
the system descriptor $\xi$, whereas spatial, mass-equivalent, and
temporal resolutions belong to an observation protocol $\lambda$.
Their anchor-independent comparison is the resolution margin
$g_i=\log_b(q_i/\varepsilon_i)$. Monomial physical relations become
linear constraint surfaces in log coordinates. In Planck normalization,
the causal condition $L\leq cT$ gives $\chi_L\leq\chi_T$, the reduced
Compton relation gives $\chi_L=-\chi_M$, and the Schwarzschild relation
gives $\chi_L=\chi_M+\log_b2$. A reproducible cosmological ledger then
shows why a bare ``mass of the observable Universe'' is not a typed
quantity: baryonic, total-matter, and total-energy mass equivalents
differ by more than a decade under the same rounded horizon convention.
\end{abstract}

\tableofcontents

\section{Scope and claim boundary}

The original ruler note correctly normalized physical quantities before
taking logarithms. Its unresolved issue was semantic rather than
algebraic: object sizes, instrument resolutions, and cosmological
landmarks were placed in the same coordinates without distinct
passports. The notation
\[
  Y_X(s,\mu,\tau)=O_{s,\mu,\tau}(X)
\]
also treated every point of $\R^3$ as if it automatically specified an
observation channel. It does not. A channel requires a declared
measurement model, frame, calibration, noise law, estimator, and
applicable axes.

\begin{warning}[Claim firewall]
The Planck length and Planck time are not asserted to be experimentally
proved minimum scales. The present comoving particle horizon is not the
diameter of the entire Universe. Similar numerical spans of roughly
sixty decades do not imply a new physical law. A camera, clock, or
spectrometer need not measure all of length, invariant mass, and time.
\end{warning}

The strict claim is limited to the following:
\begin{quote}
Positive physical scales admit useful typed logarithmic coordinates.
Systems generally determine labelled sets or regions of such
coordinates, while observation protocols determine separate resolution
coordinates and response traces.
\end{quote}

\section{Typed scale data and the two-passport rule}

\begin{definition}[Scale item]
A scale item is a record
\[
 \sigma=(\mathsf{id},\mathsf{axis},q,\mathsf{role},
         \mathsf{frame},\mathsf{convention},\mathsf U),
\]
where $\mathsf{axis}\in\{L,M,T\}$, $q>0$ has the corresponding physical
dimension, $\mathsf{role}$ identifies the measurand, and $\mathsf U$
contains exactness or uncertainty information.
\end{definition}

The role is indispensable. A particle radius, a wavelength, a correlation
length, and a detector pixel pitch all have dimension $L$, but they are
not interchangeable. Likewise, invariant rest mass is not relativistic
energy divided by an undeclared constant, and proper time is not a
coordinate duration without a clock congruence.

\begin{definition}[Intrinsic scale catalogue]
For a system $X$, let
\[
 \Xi_X=\{\sigma_\alpha:\alpha\in A_X\}
\]
be a labelled, possibly partial and multivalued catalogue of intrinsic
or model-derived scales. It is a component of the system descriptor
$\xi_X$. A missing axis is recorded as \texttt{not-applicable}; it is
not encoded by $0$ or $+\infty$.
\end{definition}

\begin{definition}[Observation protocol passport]
An observation protocol is a tuple
\[
\lambda=
\bigl(
R_\lambda,O_\lambda,K_\lambda,Q_\lambda,
\bm\varepsilon_\lambda,\bm r_\lambda^*,
\mathsf F_\lambda,\widehat\theta_\lambda,
\mathsf U_\lambda
\bigr).
\]
Here $R_\lambda$ is a reduction, $O_\lambda$ is an ideal deterministic
observation, $K_\lambda$ is a stochastic kernel or an explicit noiseless
flag, $Q_\lambda$ is a quantizer or partition,
$\bm\varepsilon_\lambda$ contains only applicable resolutions,
$\bm r_\lambda^*$ contains their reference scales,
$\mathsf F_\lambda$ records frames and conventions,
$\widehat\theta_\lambda$ is an estimator or declared direct readout, and
$\mathsf U_\lambda$ records uncertainty.
\end{definition}

\begin{center}
\begin{tikzpicture}[
  node distance=11mm and 16mm,
  box/.style={draw,rounded corners,minimum width=34mm,minimum height=11mm,align=center},
  arr/.style={-{Latex[length=2.1mm]},thick}
]
\node[box,fill=lightblue] (xi) {$\xi_X,\ \Xi_X$\\system scales};
\node[box,fill=lightred,right=of xi] (lambda) {$\lambda$\\resolutions and channel};
\node[box,fill=lightgreen,below=of $(xi)!0.5!(lambda)$] (trace)
  {$\Tcal_X$\\protocol-indexed data law};
\draw[arr] (xi) -- (trace);
\draw[arr] (lambda) -- (trace);
\end{tikzpicture}
\end{center}

\begin{remark}
A mass axis in $\lambda$ exists only if the channel measures mass,
invariant mass, or a declared mass-equivalent. If the actual threshold is
an energy $\varepsilon_E$ and the model uses $E=Mc^2$, then
\[
 \varepsilon_M=\frac{\varepsilon_E}{c^2}
\]
is a typed converter. Without such a converter, ``mass resolution'' is
not available.
\end{remark}

\section{The logarithmic scale chart}

\begin{definition}[Positive scale cone and chart]
Let
\[
\Scal_{LMT}:=
\R_{>0}^{(L)}\times\R_{>0}^{(M)}\times\R_{>0}^{(T)}.
\]
For a positive reference triple
$\bm a=(a_L,a_M,a_T)$ and $b>1$, define
\[
\bm\chi_{b,\bm a}(L,M,T)
=
\begin{pmatrix}
\chiL\\ \chiM\\ \chiT
\end{pmatrix}
:=
\begin{pmatrix}
\log_b(L/a_L)\\
\log_b(M/a_M)\\
\log_b(T/a_T)
\end{pmatrix}.
\]
\end{definition}

\begin{theorem}[Global chart]
\label{thm:global-chart}
The map $\bm\chi_{b,\bm a}:\Scal_{LMT}\to\R^3$ is a smooth
diffeomorphism with inverse
\[
 \bm\chi_{b,\bm a}^{-1}(x_L,x_M,x_T)
 =
 \bigl(a_Lb^{x_L},a_Mb^{x_M},a_Tb^{x_T}\bigr).
\]
\end{theorem}

\begin{proof}
Each one-dimensional map $q\mapsto\log_b(q/a)$ is smooth, strictly
increasing, and bijective from $\Rpos$ to $\R$, with the displayed smooth
inverse. The product map has the same properties.
\end{proof}

The chart is an exact reparameterization, not a reduction. Information
loss enters when a scale item is selected from a system, rounded,
thresholded, or passed through an observation channel.

\begin{definition}[Scale region]
For a catalogue $\Xi_X$, a declared extraction rule
$E_{\mathsf r}$ selects compatible triples or partial tuples. Its
logarithmic image is
\[
 \mathscr R_{X,\mathsf r}
 =
 \bm\chi_{b,\bm a}\bigl(E_{\mathsf r}(\Xi_X)\bigr).
\]
An extended or multiscale object is represented by this labelled set or
region, not by a forced single point.
\end{definition}

\begin{warning}[No canonical triple]
An atom has several length scales and several characteristic times. A
galaxy has luminous, half-mass, virial, and halo radii. Selecting one
value on each axis is a protocol or modelling decision. The phrase
``the object occupies the point $(\chi_L,\chi_M,\chi_T)$'' is valid only
after the extraction rule has been declared.
\end{warning}

\section{Unit, anchor, and base transformation laws}

\begin{theorem}[Coordinate covariance]
\label{thm:coordinate-covariance}
Let $q,a>0$ have the same physical dimension.
\begin{enumerate}[label=(\roman*)]
\item Under a coherent unit conversion
      $(q,a)\mapsto(\alpha q,\alpha a)$ with $\alpha>0$,
      $\log_b(q/a)$ is unchanged.
\item If the reference changes from $a$ to $a'>0$, then
      \[
      \chi_{b,a'}(q)
      =
      \chi_{b,a}(q)-\log_b\frac{a'}a.
      \]
\item If the base changes from $b$ to $b'>1$, then
      \[
      \chi_{b',a}(q)
      =
      \frac{\ln b}{\ln b'}\,\chi_{b,a}(q).
      \]
\end{enumerate}
\end{theorem}

\begin{proof}
The first statement follows from cancellation of $\alpha$. The second is
the logarithm of $(q/a)/(a'/a)$. The third follows from the
change-of-base formula.
\end{proof}

Thus an anchor change translates log-space. Coordinate differences,
axis orderings, and any construction depending only on differences are
anchor independent. A change from decades to octaves or nats rescales
the chart.

\begin{definition}[Weighted log-scale distance]
For a symmetric positive-definite dimensionless matrix
$W\in\R^{3\times3}$, define
\[
d_{W,b}(\bm q,\bm q')
:=
\sqrt{
\bigl(\bm\chi_{b,\bm a}(\bm q)
      -\bm\chi_{b,\bm a}(\bm q')\bigr)^\mathsf T
W
\bigl(\bm\chi_{b,\bm a}(\bm q)
      -\bm\chi_{b,\bm a}(\bm q')\bigr)
}.
\]
\end{definition}

\begin{theorem}[Metric and transformation law]
\label{thm:log-metric}
The function $d_{W,b}$ is a metric on $\Scal_{LMT}$, independent of the
anchor $\bm a$ and invariant under coherent unit changes. Under a base
change $b\mapsto b'$, it remains numerically invariant if
\[
 W_{b'}
 =
 \left(\frac{\ln b'}{\ln b}\right)^2W_b.
\]
\end{theorem}

\begin{proof}
The chart in \cref{thm:global-chart} is injective, and the
positive-definite quadratic form defines a norm on coordinate
differences. Anchor translations cancel in those differences. Unit
invariance follows from \cref{thm:coordinate-covariance}. If
$c_{bb'}=\ln b/\ln b'$, then $\Delta\bm\chi_{b'}=c_{bb'}\Delta\bm\chi_b$;
the displayed transformation of $W$ cancels $c_{bb'}^2$.
\end{proof}

\begin{warning}[Metric convention]
$W=I$ declares that one decade of length, mass, and time contributes
equally. This may be useful for visualization, but it is not a physical
law. A scientifically interpretable $W$ must come from a decision
problem, tolerances, or a covariance model.
\end{warning}

\section{Descriptor coordinates versus resolution coordinates}

The Distance-Scale Interface uses a resolution coordinate whose
orientation is opposite to the descriptor coordinate.

\begin{definition}[Protocol resolution depth]
For an applicable protocol axis $i$ with resolution
$\varepsilon_i>0$ and reference $r_i^*>0$, define
\[
 \eta_i(\lambda)
 :=
 \log_b\frac{r_i^*}{\varepsilon_i}.
\]
Finer resolution increases $\eta_i$. If $r_i^*=a_i$, then
\[
 \eta_i=-\chi_{b,a_i}(\varepsilon_i).
\]
\end{definition}

\begin{definition}[Resolution margin]
For a system scale $q_i>0$ and protocol resolution
$\varepsilon_i>0$ on the same typed axis, define
\[
 g_i(q_i,\lambda)
 :=
 \log_b\frac{q_i}{\varepsilon_i}.
\]
\end{definition}

\begin{proposition}[Anchor-independent comparison]
\label{prop:margin}
For any common positive anchor $a_i$,
\[
g_i
=
\chi_{b,a_i}(q_i)-\chi_{b,a_i}(\varepsilon_i).
\]
Hence $g_i$ is invariant under coherent unit changes and common anchor
changes.
\end{proposition}

\begin{proof}
Subtract the two logarithms and combine their arguments.
\end{proof}

In a simple scale-separation model, $g_i>0$ means that the intrinsic
scale exceeds the nominal resolution. It does not by itself specify
contrast, signal-to-noise ratio, inverse-problem conditioning, or a
detection probability.

\begin{definition}[Protocol-indexed observation trace]
For a declared family $\Lambda$ of protocols, a deterministic trace is
\[
\Tcal_X^{\mathrm{det}}
=
\set{
\bigl(\bm\eta(\lambda),O_\lambda(X)\bigr):
\lambda\in\Lambda
}.
\]
For a stochastic channel it is
\[
\Tcal_X^{\mathrm{stoch}}
=
\set{
\bigl(\bm\eta(\lambda),
\Law(Y_\lambda\mid X,\lambda)\bigr):
\lambda\in\Lambda
}.
\]
\end{definition}

The shorthand $O_{\bm\eta}$ is admissible only after the map
$\bm\eta\mapsto\lambda(\bm\eta)$ and all nuisance parameters have been
defined.

\section{Monomial algebra and constraint surfaces}

\begin{theorem}[Log-linearization of positive monomial laws]
\label{thm:monomial}
Let positive inputs $\bm q=(q_1,\ldots,q_n)$ and positive outputs
$r_i$ satisfy
\[
r_i=C_i\prod_{j=1}^n q_j^{A_{ij}},
\qquad
r_i^*=C_i\prod_{j=1}^n(q_j^*)^{A_{ij}},
\]
where the references are dimensionally matched. Then
\[
\log_b\frac{r_i}{r_i^*}
=
\sum_{j=1}^n A_{ij}
\log_b\frac{q_j}{q_j^*}.
\]
In vector form, $\bm\chi_r=A\bm\chi_q$.
\end{theorem}

\begin{proof}
Divide the two monomials, cancel $C_i$, and apply the logarithm product
and power laws.
\end{proof}

This theorem is the precise structural reason that dimensional
monomials become hyperplanes in logarithmic coordinates. It is a
coordinate simplification, not a claim that arbitrary nonlinear physics
becomes linear.

\subsection{Planck identities}

Use the standard definitions
\[
\ell_P=\sqrt{\frac{\hbar G}{c^3}},
\qquad
m_P=\sqrt{\frac{\hbar c}{G}},
\qquad
t_P=\sqrt{\frac{\hbar G}{c^5}}.
\]
Direct substitution gives
\[
t_P=\frac{\ell_P}{c},
\qquad
\ell_Pm_P=\frac{\hbar}{c},
\qquad
m_Pc^2=\frac{\hbar}{t_P}.
\]
The three Planck anchors are therefore built from common constants and
are not statistically or conceptually independent axes.

\subsection{Causal, quantum, and gravitational examples}

\begin{corollary}[Causal half-space]
\label{cor:causal}
If $L$ and $T$ refer to the same causal process and $L\leq cT$, then in
Planck normalization
\[
\chiL\leq\chiT.
\]
\end{corollary}

\begin{proof}
Because $\ell_P=ct_P$,
\[
\frac{L/\ell_P}{T/t_P}=\frac{L}{cT}\leq1.
\]
Apply the increasing logarithm.
\end{proof}

If $L$ and $T$ belong to unrelated catalogue entries, the hypothesis of
\cref{cor:causal} fails and no such inequality follows.

\begin{corollary}[Mass-energy, density, Compton, and Schwarzschild planes]
\label{cor:planes}
In Planck normalization:
\begin{align*}
E=Mc^2,\quad E_P=m_Pc^2
&\Longrightarrow
\chi_E=\chiM,\\
\rho=M/L^3,\quad \rho_P=m_P/\ell_P^3
&\Longrightarrow
\chi_\rho=\chiM-3\chiL,\\
\bar\lambda_C=\hbar/(Mc)
&\Longrightarrow
\chi_{\bar\lambda_C}=-\chiM,\\
r_s=2GM/c^2
&\Longrightarrow
\chi_{r_s}=\chiM+\log_b2.
\end{align*}
\end{corollary}

\begin{proof}
The first two identities follow from \cref{thm:monomial}. For the
reduced Compton length, use $\ell_Pm_P=\hbar/c$:
\[
\frac{\bar\lambda_C}{\ell_P}=\frac{m_P}{M}.
\]
For the Schwarzschild radius, $Gm_P/c^2=\ell_P$, hence
\[
\frac{r_s}{\ell_P}=2\frac{M}{m_P}.
\]
Taking logarithms proves the result.
\end{proof}

\begin{warning}[Physical-domain boundary]
The lines in \cref{cor:planes} are relations among specifically coupled
measurands. They do not turn the whole scale catalogue into one
three-dimensional state space, and they do not by themselves establish
a quantum-gravity model.
\end{warning}

\section{Planck anchors and coordinate uncertainty}

{\small
\begin{longtable}{L{0.17\linewidth}L{0.25\linewidth}L{0.22\linewidth}L{0.20\linewidth}}
\toprule
Anchor & Definition & 2022 CODATA value & Status\\
\midrule
Planck length $\ell_P$ & $\sqrt{\hbar G/c^3}$ & $1.61625\times 10^{-35}\,\mathrm{m}$ & normalization, not a proved minimum length\\
Planck mass $m_P$ & $\sqrt{\hbar c/G}$ & $2.17643\times 10^{-8}\,\mathrm{kg}$ & normalization, not a minimum mass\\
Planck time $t_P$ & $\sqrt{\hbar G/c^5}$ & $5.39125\times 10^{-44}\,\mathrm{s}$ & normalization, not a proved minimum duration\\
\bottomrule
\end{longtable}
}


The 2022 CODATA values inherit uncertainty primarily through $G$; $c$
and $h$ are exact defining constants in the revised SI. Consequently,
the Planck anchors are correlated. Treating their listed standard
uncertainties as independent can give an incorrect covariance for
derived scale coordinates.

\begin{proposition}[First-order uncertainty of a log ratio]
\label{prop:uncertainty}
Let
\[
\chi=\log_b(q/a)
\]
with $q,a>0$. To first order,
\[
u^2(\chi)
=
\frac{1}{(\ln b)^2}
\left[
\frac{u^2(q)}{q^2}
+\frac{u^2(a)}{a^2}
-2\frac{\Cov(q,a)}{qa}
\right].
\]
For several correlated inputs,
\[
\Sigma_\chi=J\Sigma_{q,a}J^\mathsf T,
\]
where $J$ is the Jacobian of the logarithmic map.
\end{proposition}

\begin{proof}
Differentiate
$\chi=(\ln q-\ln a)/\ln b$ and apply the first-order covariance
propagation formula.
\end{proof}

The formula is a local approximation. Large, asymmetric, or
model-posterior uncertainties should be propagated by the full
distribution rather than by a symmetric linearized error bar.

\section{Cosmological landmarks are typed model outputs}

The phrase ``observable-universe diameter'' must name a cosmological
distance. The legacy $93$ Gly entry is retained only as a rounded
\emph{present comoving particle-horizon diameter} in a declared
$a_0=1$ convention. It is not a light-travel distance and not the size
of the entire Universe. Cosmological distance definitions are distinct
even when they share the SI dimension of length.

The mass axis has a larger ambiguity. For a flat-volume illustrative
ledger, let
\[
\rho_{c0}=\frac{3H_0^2}{8\pi G},
\qquad
V_{\mathrm{flat}}=\frac{4\pi}{3}R_{\mathrm{ph},0}^3.
\]
Then three different mass-equivalent landmarks are
\[
M_b=\Omega_b\rho_{c0}V_{\mathrm{flat}},
\qquad
M_m=\Omega_m\rho_{c0}V_{\mathrm{flat}},
\qquad
M_E=\rho_{c0}V_{\mathrm{flat}}.
\]
$M_b$ counts the baryonic component, $M_m$ all pressureless matter in
the adopted parameterization, and $M_E$ the total critical
mass-equivalent. Dark energy is not a collection of rest-mass particles;
the last quantity is explicitly an energy-equivalent bookkeeping scale.

{\small
\begin{longtable}{L{0.25\linewidth}L{0.20\linewidth}L{0.11\linewidth}L{0.25\linewidth}}
\toprule
Landmark & Value in SI & Coordinate & Convention\\
\midrule
Present comoving particle-horizon diameter & $8.79848\times 10^{26}\,\mathrm{m}$ & $61.7359$ & rounded 93 Gly landmark\\
FLRW cosmic age & $4.35400\times 10^{17}\,\mathrm{s}$ & $60.9072$ & Planck 2018 baseline, $13.797\pm0.023$ Gyr\\
Baryonic mass-equivalent & $1.50025\times 10^{53}\,\mathrm{kg}$ & $60.8384$ & $\Omega_b\rho_{c0}V$\\
All-matter mass-equivalent & $9.58576\times 10^{53}\,\mathrm{kg}$ & $61.6439$ & $\Omega_m\rho_{c0}V$\\
Total-energy mass-equivalent & $3.04310\times 10^{54}\,\mathrm{kg}$ & $62.1456$ & $\rho_{c0}V$\\
\bottomrule
\end{longtable}
}


The table uses a rounded $93$ Gly diameter and a Planck 2018 flat
base-$\Lambda$CDM parameter ledger:
$H_0=67.4\,\mathrm{km\,s^{-1}\,Mpc^{-1}}$,
$\Omega_m=0.315$, and $\Omega_b=0.0493$. Because the rounded horizon
entry and the full parameter covariance are not supplied, no combined
confidence interval is fabricated. The values are reproducible
landmarks, not a precision cosmological analysis.

\begin{proposition}[The legacy mass anchor is ambiguous]
Under this single illustrative ledger,
\[
M_b\approx1.50\times10^{53}\,\mathrm{kg},\qquad
M_m\approx9.59\times10^{53}\,\mathrm{kg},\qquad
M_E\approx3.04\times10^{54}\,\mathrm{kg}.
\]
Their Planck-mass coordinates are approximately
$60.84$, $61.64$, and $62.15$, respectively. Therefore an unqualified
$10^{53}\,\mathrm{kg}$ entry cannot serve as a unique cosmic mass
anchor.
\end{proposition}

\begin{proof}
Substitute the declared $H_0$, density fractions, and flat-volume
convention into the preceding definitions, then divide by $m_P$ and take
$\log_{10}$.
\end{proof}

\begin{warning}[Epoch and model dependence]
Changing the cosmological model, horizon definition, spatial curvature,
epoch, or matter component changes these landmarks. They are catalogue
records, not constants defining the logarithmic coordinate system.
\end{warning}

\section{Selected reproducible catalogue}

The following table is generated from the machine-readable input ledger.
Entries marked ``order of magnitude'' or ``declared reference'' are
visual landmarks rather than measurements. Every coordinate uses the
corresponding 2022 CODATA Planck anchor and base $10$.

{\small
\begin{longtable}{L{0.21\linewidth}L{0.08\linewidth}L{0.20\linewidth}L{0.12\linewidth}L{0.21\linewidth}}
\toprule
Entry & Axis & SI value & Coordinate & Status\\
\midrule
\endhead
planck length & length & $1.6163\times 10^{-35}\,\mathrm{m}$ & $0.000$ & CODATA recommended\\
atomic order & length & $1.0000\times 10^{-10}\,\mathrm{m}$ & $24.791$ & order of magnitude marker\\
human length reference & length & $1.0000\times 10^{0}\,\mathrm{m}$ & $34.791$ & declared not population estimate\\
earth diameter & length & $1.2742\times 10^{7}\,\mathrm{m}$ & $41.897$ & rounded mean diameter\\
astronomical unit & length & $1.4960\times 10^{11}\,\mathrm{m}$ & $45.966$ & exact IAU definition\\
observable particle horizon diameter & length & $8.7985\times 10^{26}\,\mathrm{m}$ & $61.736$ & rounded model dependent\\
electron mass & mass & $9.1094\times 10^{-31}\,\mathrm{kg}$ & $-22.378$ & CODATA recommended\\
proton mass & mass & $1.6726\times 10^{-27}\,\mathrm{kg}$ & $-19.114$ & CODATA recommended\\
planck mass & mass & $2.1764\times 10^{-8}\,\mathrm{kg}$ & $0.000$ & CODATA recommended\\
human mass reference & mass & $7.0000\times 10^{1}\,\mathrm{kg}$ & $9.507$ & declared not population estimate\\
earth mass & mass & $5.9722\times 10^{24}\,\mathrm{kg}$ & $32.438$ & rounded astronomical value\\
solar mass & mass & $1.9885\times 10^{30}\,\mathrm{kg}$ & $37.961$ & rounded astronomical value\\
planck time & time & $5.3912\times 10^{-44}\,\mathrm{s}$ & $0.000$ & CODATA recommended\\
femtosecond & time & $1.0000\times 10^{-15}\,\mathrm{s}$ & $28.268$ & SI prefix exact\\
second & time & $1.0000\times 10^{0}\,\mathrm{s}$ & $43.268$ & SI base unit\\
day & time & $8.6400\times 10^{4}\,\mathrm{s}$ & $48.205$ & conventional exact duration\\
Julian year & time & $3.1558\times 10^{7}\,\mathrm{s}$ & $50.767$ & conventional exact duration\\
universe age & time & $4.3540\times 10^{17}\,\mathrm{s}$ & $60.907$ & model inferred parameter\\
\bottomrule
\end{longtable}
}


The equality of the light-year length coordinate and the Julian-year
time coordinate to the displayed precision is not mysterious:
$1\,\mathrm{ly}=c\times1\,\mathrm{Julian\ year}$ and
$\ell_P=ct_P$. It is an exact consequence of matched definitions.

\section{Normative reporting contract}

For a static catalogue, stochastic-channel, quantizer, temporal
resolution, and observation-horizon fields may be
\texttt{not-applicable}; the field must still be present. For an actual
measurement trace, applicable fields require values rather than prose
placeholders.

\begin{longtable}{L{0.23\textwidth}L{0.20\textwidth}L{0.47\textwidth}}
\toprule
Field & Type & Required declaration\\
\midrule
\endhead
Hidden system $X$ & state space or object & domain and admissible states\\
System descriptor $\xi_X$ & typed record & intrinsic model data and
scale catalogue $\Xi_X$\\
Scale item & typed positive quantity & axis, role, frame, convention,
value or interval, uncertainty\\
Extraction rule & reduction & which labelled items form a tuple or
region\\
Reference triple $\bm a$ & positive typed tuple & Planck or alternative
anchors, with source\\
Logarithm base $b$ & dimensionless, $b>1$ & decades, octaves, or other
declared base\\
Observation protocol $\lambda$ & channel passport & reduction,
observation, noise/exactness, quantizer, estimator\\
Applicable axes & subset of $\{L,M,T\}$ & missing axes recorded as
not applicable\\
Spatial resolution & positive length or not applicable & frame and
calibration\\
Mass resolution & positive mass or not applicable & direct channel or
explicit energy/mass converter\\
Temporal resolution & positive duration or not applicable & clock and
frame convention\\
Observation horizon & duration or not applicable & acquisition interval\\
Reference scales $\bm r^*$ & positive typed tuple & normalization of
resolution depth\\
Resolution margin & dimensionless & matched measurand and resolution\\
Frame passport & metadata & proper, coordinate, rest, comoving, or other
declared convention\\
Cosmology passport & model metadata or not applicable & model, epoch,
distance, volume, and component definitions\\
Uncertainty & covariance, law, or exact flag & dependence and propagation
semantics\\
Weight matrix $W$ & positive definite or not applicable & rationale and
base-transformation policy\\
Defect normalization & typed ratio or not applicable & required before
scalar loss aggregation\\
\bottomrule
\end{longtable}

\section{Claim audit}

\begin{longtable}{L{0.43\textwidth}L{0.19\textwidth}L{0.28\textwidth}}
\toprule
Statement & Status & Necessary hypotheses\\
\midrule
\endhead
$\bm\chi_{b,\bm a}$ is a global chart on $(\Rpos)^3$ & theorem &
positive scales, positive anchors, $b>1$\\
Log coordinates are invariant under a unit change & theorem & quantity
and anchor converted coherently\\
Anchor changes preserve coordinate differences & theorem & common typed
anchor change\\
$d_{W,b}$ is a scale-space metric & theorem & $W$ symmetric positive
definite\\
$W=I$ is the unique physical metric & false & no such canonical choice
is proved\\
Every system has one canonical $(L,M,T)$ point & false & requires a
declared extraction rule\\
Every protocol measures all three axes & false & missing axes are
allowed\\
$g_i>0$ guarantees detection & false & only a simple scale-separation
diagnostic\\
Planck units are minimum physical scales & not claimed & they are used as
normalizations\\
The causal inequality gives $\chi_L\leq\chi_T$ & corollary & paired
length and time of the same causal process\\
Compton and Schwarzschild relations become affine lines & corollary &
declared standard formulas and Planck normalization\\
$93$ Gly is the diameter of the entire Universe & false & rounded present
comoving particle-horizon convention only\\
$10^{53}\,\mathrm{kg}$ is the unique mass of the observable Universe &
false & component, volume, epoch, and model are required\\
Three spans near sixty decades imply a new law & not claimed & numerical
similarity depends on selected landmarks\\
\bottomrule
\end{longtable}

\section{Conclusion}

The strict Planck-to-cosmos ruler object is not a bare point in $\R^3$.
It is the pair
\[
\bigl(
\xi_X,\Xi_X;
\lambda,\bm\varepsilon_\lambda,\bm r_\lambda^*
\bigr),
\]
together with a positive reference tuple, logarithm base, frame and
cosmology conventions, uncertainty model, and a protocol-indexed
observation trace. The descriptor chart
\[
\bm\chi_{b,\bm a}(\bm q)
=
\log_b(\bm q/\bm a)
\]
is unit invariant and turns positive monomial relations into linear
constraints. The resolution margin
$\log_b(q_i/\varepsilon_i)$ compares a typed system scale with a typed
protocol resolution without conflating the two. Cosmic quantities are
optional model outputs. This separation closes the central blocker in
v1.0 and makes the module compatible with the GO Core distance-scale
interface.

\begin{thebibliography}{12}

\bibitem{CODATA2022}
P. J. Mohr, D. B. Newell, B. N. Taylor, and E. Tiesinga,
``CODATA Recommended Values of the Fundamental Physical Constants:
2022,''
\emph{Reviews of Modern Physics} 97 (2025), 025002.

\bibitem{NIST}
National Institute of Standards and Technology,
\emph{Reference on Constants, Units, and Uncertainty: 2022 CODATA
recommended values}, accessed 2026.

\bibitem{BIPM}
Bureau International des Poids et Mesures,
\emph{The International System of Units (SI Brochure)},
9th ed., version 4.01, 2026.

\bibitem{GUM}
Joint Committee for Guides in Metrology,
\emph{Evaluation of Measurement Data --- Guide to the Expression of
Uncertainty in Measurement}, JCGM 100:2008.

\bibitem{PlanckOverview}
Planck Collaboration,
``Planck 2018 results. I. Overview and the cosmological legacy of
Planck,''
\emph{Astronomy \& Astrophysics} 641 (2020), A1.

\bibitem{PlanckVI}
Planck Collaboration,
``Planck 2018 results. VI. Cosmological parameters,''
\emph{Astronomy \& Astrophysics} 641 (2020), A6.

\bibitem{Hogg}
D. W. Hogg,
``Distance measures in cosmology,''
arXiv:astro-ph/9905116 (1999).

\bibitem{BBI}
D. Burago, Y. Burago, and S. Ivanov,
\emph{A Course in Metric Geometry},
American Mathematical Society, 2001.

\end{thebibliography}

\end{document}
